\IfFileExists{stacks-project.cls}{%
\documentclass{stacks-project}
}{%
\documentclass{amsart}
}
\usepackage{amssymb}

% For dealing with references we use the comment environment
\usepackage{verbatim}
\usepackage{fancyvrb}
\newenvironment{reference}{\comment}{\endcomment}
%\newenvironment{reference}{}{}
\newenvironment{slogan}{\comment}{\endcomment}
\newenvironment{history}{\comment}{\endcomment}

% For source figures retained as assets
\usepackage{graphicx}

% For commutative diagrams we use Xy-pic
\usepackage[all]{xy}

% We use 2cell for 2-commutative diagrams.
\xyoption{2cell}
\UseAllTwocells

% We use multicol for the list of chapters between chapters
\usepackage{multicol}

% This is generally recommended for better output
\usepackage{lmodern}
\usepackage[T1]{fontenc}
\usepackage{textalpha}

% For cross-file-references
% We prefer the snapshot tmp/xref/#2.aux,
% which latexmkrc replaces atomically, so
% that documents compiled simultaneously
% never read a half-written .aux file.
% The last argument keeps links to #2.pdf.
\usepackage{xr-hyper}
\newcommand{\maybeexternaldocument}[2]{%
  \IfFileExists{#2.aux}{%
    \externaldocument[#1]{#2}%
  }{%
    \PackageWarning{stack}{Missing #2.aux; external references with prefix #1 unavailable}%
  }%
}
\let\externaldocumentorig\externaldocument
\renewcommand{\externaldocument}[2][]{%
  \IfFileExists{tmp/xref/#2.aux}{%
    \externaldocumentorig[#1]{tmp/xref/#2}[#2.pdf]%
  }{%
  \IfFileExists{#2.aux}{%
    \externaldocumentorig[#1]{#2}%
  }{%
    \PackageWarning{stack}{Missing #2.aux; external references with prefix #1 unavailable}%
  }%
  }%
}

% Package for hypertext links:
\usepackage{hyperref}

% For any local file, say "hello.tex" you want to link to please
% use \externaldocument[hello-]{hello}
\externaldocument[introduction-]{introduction}
\externaldocument[conventions-]{conventions}
\externaldocument[sets-]{sets}
\externaldocument[categories-]{categories}
\externaldocument[topology-]{topology}
\externaldocument[sheaves-]{sheaves}
\externaldocument[sites-]{sites}
\externaldocument[stacks-]{stacks}
\externaldocument[fields-]{fields}
\externaldocument[algebra-]{algebra}
\externaldocument[brauer-]{brauer}
\externaldocument[homology-]{homology}
\externaldocument[derived-]{derived}
\externaldocument[simplicial-]{simplicial}
\externaldocument[more-algebra-]{more-algebra}
\externaldocument[smoothing-]{smoothing}
\externaldocument[modules-]{modules}
\externaldocument[sites-modules-]{sites-modules}
\externaldocument[injectives-]{injectives}
\externaldocument[cohomology-]{cohomology}
\externaldocument[sites-cohomology-]{sites-cohomology}
\externaldocument[dga-]{dga}
\externaldocument[dpa-]{dpa}
\externaldocument[sdga-]{sdga}
\externaldocument[hypercovering-]{hypercovering}
\externaldocument[schemes-]{schemes}
\externaldocument[constructions-]{constructions}
\externaldocument[properties-]{properties}
\externaldocument[morphisms-]{morphisms}
\externaldocument[coherent-]{coherent}
\externaldocument[divisors-]{divisors}
\externaldocument[limits-]{limits}
\externaldocument[varieties-]{varieties}
\externaldocument[topologies-]{topologies}
\externaldocument[descent-]{descent}
\externaldocument[perfect-]{perfect}
\externaldocument[more-morphisms-]{more-morphisms}
\externaldocument[flat-]{flat}
\externaldocument[groupoids-]{groupoids}
\externaldocument[more-groupoids-]{more-groupoids}
\externaldocument[etale-]{etale}
\externaldocument[chow-]{chow}
\externaldocument[intersection-]{intersection}
\externaldocument[pic-]{pic}
\externaldocument[weil-]{weil}
\externaldocument[adequate-]{adequate}
\externaldocument[dualizing-]{dualizing}
\externaldocument[duality-]{duality}
\externaldocument[discriminant-]{discriminant}
\externaldocument[derham-]{derham}
\externaldocument[local-cohomology-]{local-cohomology}
\externaldocument[algebraization-]{algebraization}
\externaldocument[curves-]{curves}
\externaldocument[resolve-]{resolve}
\externaldocument[models-]{models}
\externaldocument[functors-]{functors}
\externaldocument[equiv-]{equiv}
\externaldocument[pione-]{pione}
\externaldocument[etale-cohomology-]{etale-cohomology}
\externaldocument[proetale-]{proetale}
\externaldocument[relative-cycles-]{relative-cycles}
\externaldocument[more-etale-]{more-etale}
\externaldocument[trace-]{trace}
\externaldocument[crystalline-]{crystalline}
\externaldocument[spaces-]{spaces}
\externaldocument[spaces-properties-]{spaces-properties}
\externaldocument[spaces-morphisms-]{spaces-morphisms}
\externaldocument[decent-spaces-]{decent-spaces}
\externaldocument[spaces-cohomology-]{spaces-cohomology}
\externaldocument[spaces-limits-]{spaces-limits}
\externaldocument[spaces-divisors-]{spaces-divisors}
\externaldocument[spaces-over-fields-]{spaces-over-fields}
\externaldocument[spaces-topologies-]{spaces-topologies}
\externaldocument[spaces-descent-]{spaces-descent}
\externaldocument[spaces-perfect-]{spaces-perfect}
\externaldocument[spaces-more-morphisms-]{spaces-more-morphisms}
\externaldocument[spaces-flat-]{spaces-flat}
\externaldocument[spaces-groupoids-]{spaces-groupoids}
\externaldocument[spaces-more-groupoids-]{spaces-more-groupoids}
\externaldocument[bootstrap-]{bootstrap}
\externaldocument[spaces-pushouts-]{spaces-pushouts}
\externaldocument[spaces-chow-]{spaces-chow}
\externaldocument[groupoids-quotients-]{groupoids-quotients}
\externaldocument[spaces-more-cohomology-]{spaces-more-cohomology}
\externaldocument[spaces-simplicial-]{spaces-simplicial}
\externaldocument[spaces-duality-]{spaces-duality}
\externaldocument[formal-spaces-]{formal-spaces}
\externaldocument[restricted-]{restricted}
\externaldocument[spaces-resolve-]{spaces-resolve}
\externaldocument[formal-defos-]{formal-defos}
\externaldocument[defos-]{defos}
\externaldocument[cotangent-]{cotangent}
\externaldocument[examples-defos-]{examples-defos}
\externaldocument[algebraic-]{algebraic}
\externaldocument[examples-stacks-]{examples-stacks}
\externaldocument[stacks-sheaves-]{stacks-sheaves}
\externaldocument[criteria-]{criteria}
\externaldocument[artin-]{artin}
\externaldocument[quot-]{quot}
\externaldocument[stacks-properties-]{stacks-properties}
\externaldocument[stacks-morphisms-]{stacks-morphisms}
\externaldocument[stacks-limits-]{stacks-limits}
\externaldocument[stacks-cohomology-]{stacks-cohomology}
\externaldocument[stacks-perfect-]{stacks-perfect}
\externaldocument[stacks-introduction-]{stacks-introduction}
\externaldocument[stacks-more-morphisms-]{stacks-more-morphisms}
\externaldocument[stacks-geometry-]{stacks-geometry}
\externaldocument[moduli-]{moduli}
\externaldocument[moduli-curves-]{moduli-curves}
\externaldocument[examples-]{examples}
\externaldocument[exercises-]{exercises}
\externaldocument[guide-]{guide}
\externaldocument[desirables-]{desirables}
\externaldocument[coding-]{coding}
\externaldocument[obsolete-]{obsolete}
\externaldocument[fdl-]{fdl}
\externaldocument[index-]{index}

% Heap Project documents
\externaldocument[k3-]{k3}
\externaldocument[torelli-]{torelli}
\externaldocument[linear-algebra-]{linear-algebra}
\externaldocument[groups-]{groups}
\externaldocument[complex-analysis-]{complex-analysis}
\externaldocument[functional-analysis-]{functional-analysis}
\externaldocument[categories-]{categories}
\externaldocument[commutative-algebra-]{commutative-algebra}
\externaldocument[ringed-spaces-]{ringed-spaces}
\externaldocument[homological-algebra-]{homological-algebra}
\externaldocument[lie-algebras-]{lie-algebras}
\externaldocument[representation-theory-]{representation-theory}
\externaldocument[smooth-manifolds-]{smooth-manifolds}
\externaldocument[differential-topology-]{differential-topology}
\externaldocument[algebraic-topology-]{algebraic-topology}
\externaldocument[differential-geometry-]{differential-geometry}
\externaldocument[algebraic-geometry-]{algebraic-geometry}
\externaldocument[symplectic-geometry-]{symplectic-geometry}
\externaldocument[complex-geometry-]{complex-geometry}
\externaldocument[hodge-theory-]{hodge-theory}
\externaldocument[%
variations-of-hodge-structures-]{%
variations-of-hodge-structures}
\externaldocument[%
homological-mirror-symmetry-]{%
homological-mirror-symmetry}
\maybeexternaldocument{deformations-}{deformations}
\externaldocument[vertex-algebras-]{vertex-algebras}
\externaldocument[springer-]{springer}
\externaldocument[infty-categories-]{infty-categories}
\externaldocument[seminars-differential-geometry-]{%
seminars-differential-geometry}
\externaldocument[seminars-complex-geometry-]{%
seminars-complex-geometry}
\externaldocument[seminars-algebraic-geometry-]{%
seminars-algebraic-geometry}
\externaldocument[seminars-others-]{%
seminars-others}
\externaldocument[%
geometric-prequantization-]{%
geometric-prequantization}
\externaldocument[surfaces-]{surfaces}
\externaldocument[birational-maps-commutative-algebra-]{birational-maps-commutative-algebra}
\externaldocument[cremona-transformations-]{cremona-transformations}
\externaldocument[derived-categories-of-sheaves-]{derived-categories-of-sheaves}
\externaldocument[geometric-invariant-theory-]{geometric-invariant-theory}
\externaldocument[geometric-stability-conditions-and-group-actions-]{geometric-stability-conditions-and-group-actions}
\externaldocument[moduli-spaces-of-sheaves-]{moduli-spaces-of-sheaves}
\externaldocument[bridgeland-stability-general-theory-]{bridgeland-stability-general-theory}
\externaldocument[hilbert-schemes-]{hilbert-schemes}
\externaldocument[noncommutative-abelian-surfaces-and-kummer-type-hyperkahler-varieties-]{noncommutative-abelian-surfaces-and-kummer-type-hyperkahler-varieties}
\externaldocument[mumford-tate-groups-in-hodge-theory-]{mumford-tate-groups-in-hodge-theory}
\externaldocument[hodge-birational-atoms-2026-]{hodge-birational-atoms-2026}
\externaldocument[xxii-egd-2026-]{%
xxii-egd-2026}
\externaldocument[pde-]{pde}
\externaldocument[probability-]{probability}

\externaldocument[linguistics-]{linguistics}
\externaldocument[genealogia-familiar-]{%
genealogia-familiar}

\externaldocument[%
16th-alga-meeting-2026-]{%
16th-alga-meeting-2026}

\externaldocument[%
iv-jornadas-lefschetz-2026-]{%
iv-jornadas-lefschetz-2026}

\externaldocument[reading-the-masters-]{%
reading-the-masters}

\externaldocument[real-analysis-]{%
real-analysis}

\externaldocument[optimal-transport-]{%
optimal-transport}

\externaldocument[history-]{history}

\externaldocument[seiberg-witten-]{seiberg-witten}

\externaldocument[history-philosophy-biology-]{history-philosophy-biology}

\externaldocument[chemistry-]{chemistry}

\externaldocument[molecular-cell-biology-]{molecular-cell-biology}

\externaldocument[microbiology-]{microbiology}

\externaldocument[protistology-]{protistology}

\externaldocument[mycology-]{mycology}

\externaldocument[botany-]{botany}

\externaldocument[zoology-]{zoology}

\externaldocument[animal-histology-]{animal-histology}

\externaldocument[animal-physiology-]{animal-physiology}

\externaldocument[systematics-biogeography-]{systematics-biogeography}

\externaldocument[ecology-natural-resources-]{ecology-natural-resources}

\externaldocument[earth-sciences-]{earth-sciences}

\externaldocument[biostatistics-]{biostatistics}

\externaldocument[mathematical-systems-biology-]{mathematical-systems-biology}

\externaldocument[cimpa-el-salvador-2024-]{cimpa-el-salvador-2024}

% Theorem environments.
%
\theoremstyle{plain}
\newtheorem{theorem}[subsection]{Theorem}
\newtheorem{proposition}[subsection]{Proposition}
\newtheorem{lemma}[subsection]{Lemma}
\newtheorem{corollary}[subsection]{Corollary}

\theoremstyle{definition}
\newtheorem{definition}[subsection]{Definition}
\newtheorem{example}[subsection]{Example}
\newtheorem{exercise}[subsection]{Exercise}
\newtheorem{situation}[subsection]{Situation}

\theoremstyle{remark}
\newtheorem{remark}[subsection]{Remark}
\newtheorem{remarks}[subsection]{Remarks}

\numberwithin{equation}{subsection}

% Macros
%
\def\lim{\mathop{\mathrm{lim}}\nolimits}
\def\colim{\mathop{\mathrm{colim}}\nolimits}
\def\Spec{\mathop{\mathrm{Spec}}}
\def\Hom{\mathop{\mathrm{Hom}}\nolimits}
\def\Ext{\mathop{\mathrm{Ext}}\nolimits}
\def\SheafHom{\mathop{\mathcal{H}\!\mathit{om}}\nolimits}
\def\SheafExt{\mathop{\mathcal{E}\!\mathit{xt}}\nolimits}
\def\Sch{\mathit{Sch}}
\def\Mor{\mathop{\mathrm{Mor}}\nolimits}
\def\Ob{\mathop{\mathrm{Ob}}\nolimits}
\def\Sh{\mathop{\mathit{Sh}}\nolimits}
\def\NL{\mathop{N\!L}\nolimits}
\def\CH{\mathop{\mathrm{CH}}\nolimits}
\def\proetale{{pro\text{-}\acute{e}tale}}
\def\etale{{\acute{e}tale}}
\def\QCoh{\mathit{QCoh}}
\def\Ker{\mathop{\mathrm{Ker}}}
\def\Im{\mathop{\mathrm{Im}}}
\def\Coker{\mathop{\mathrm{Coker}}}
\def\Coim{\mathop{\mathrm{Coim}}}

% Boxtimes
%
\DeclareMathSymbol{\boxtimes}{\mathbin}{AMSa}{"02}

%
% Macros for moduli stacks/spaces
%
\def\QCohstack{\mathcal{QC}\!\mathit{oh}}
\def\Cohstack{\mathcal{C}\!\mathit{oh}}
\def\Spacesstack{\mathcal{S}\!\mathit{paces}}
\def\Quotfunctor{\mathrm{Quot}}
\def\Hilbfunctor{\mathrm{Hilb}}
\def\Curvesstack{\mathcal{C}\!\mathit{urves}}
\def\Polarizedstack{\mathcal{P}\!\mathit{olarized}}
\def\Complexesstack{\mathcal{C}\!\mathit{omplexes}}
% \Pic is the operator that assigns to X its picard group, usage \Pic(X)
% \Picardstack_{X/B} denotes the Picard stack of X over B
% \Picardfunctor_{X/B} denotes the Picard functor of X over B
\def\Pic{\mathop{\mathrm{Pic}}\nolimits}
\def\Picardstack{\mathcal{P}\!\mathit{ic}}
\def\Picardfunctor{\mathrm{Pic}}
\def\Deformationcategory{\mathcal{D}\!\mathit{ef}}


\begin{document}

\title{Mycology}
\maketitle

\phantomsection
\label{section-phantom}

\tableofcontents

\section{What this chapter covers}
\label{section-overview}

These notes cover the biology of the true
fungi: what they are, how they grow and
reproduce, how they are classified, their
fossil record, their ecology and uses,
mushroom poisoning (mycetism) and the
traditional use of mushrooms
(ethnomycology), and a collecting trip to
the pine, oak and fir forests of Hidalgo,
Mexico.

\textbf{Sources of these notes.}
The notes rest on the course \emph{Biolog\'ia
de Hongos} (third semester, Facultad de
Ciencias, UNAM, autumn 2014). No lecture
slides survive. What survives is Daniel's
own work: exam outlines on Zygomycota,
Ascomycota, Basidiomycota, Deuteromycota and
the fungus-like groups; an essay on the
evolution and systematics of fungi; a long
essay and a talk on mycetisms; a proposal
for a Lotka--Volterra model of
\emph{Sclerotinia} on sunflower; and the
team report of the field trip to Hidalgo. A
Spanish textbook chapter on Ascomycota came
with the course. The facts were checked and
completed with Webster and Weber's
textbook \cite{webster2007fungi},
Deacon's \cite{deacon2006fungal}, the
classification of Hibbett et
al.\ \cite{hibbett2007higher}, Taylor and
Berbee on divergence times
\cite{taylor2006dating}, the clinical review
of Graeme \cite{graeme2014mycetism}, the
review of \emph{Amanita muscaria} by Michelot
and Melendez-Howell
\cite{michelot2003amanita}, the psilocybin
review of Ty\v{l}s et al.\
\cite{tyls2014psilocybin}, and Hall et al.,
\emph{Edible and Poisonous Mushrooms of the
World} \cite{hall2003edible}.

\section{What fungi are}
\label{section-what-fungi-are}

About 80\,000 to 120\,000 species of fungi
have been described, out of an estimated
1.5 million
\cite[Section 1.1]{webster2007fungi}.
Fungi are heterotrophic eukaryotes that feed
by absorption: they secrete enzymes, digest
their food outside the cell, and take up the
dissolved products. Extracellular digestion
plus absorption is ``the ultimate
determinant of the fungal lifestyle''. The
typical vegetative body is a mycelium of
branching tubes (hyphae) that grow at their
tips; propagation is by numerous small
spores
\cite[Section 1.1]{webster2007fungi}.

Fungi were long classed with plants because
they do not move, absorb their food and
reproduce by spores
\cite[Section 1.5]{webster2007fungi}.
They are neither plants nor animals, and
molecular evidence places them closer to
animals than to plants
\cite[Chapter 1]{deacon2006fungal}. Like
animal sperm, the zoospores of chytrids
swim with a single posterior flagellum, a
condition called opisthokont
\cite[Section 1.4]{webster2007fungi}; the
clade of animals and fungi is named the
opisthokonts after it. The table below, after
Deacon's Table~1.1, collects the main
biochemical differences
\cite[Chapter 1]{deacon2006fungal}.

\begin{center}
\begin{tabular}{llll}
 & fungi & animals & plants \\
\hline
nutrition & absorb & ingest &
photosynthesis \\
cell wall & chitin & none &
cellulose \\
lysine & AAA & dietary &
DAP \\
sterol & ergosterol & cholesterol &
phytosterols \\
nuclei & mostly $n$ & $2n$ & $2n$ \\
\end{tabular}
\end{center}

Lysine is made by two unrelated pathways,
named after an intermediate:
$\alpha$-aminoadipic acid (AAA) and
diaminopimelic acid (DAP). The AAA pathway
occurs only in the chitin-walled fungi and
some euglenids; plants, bacteria and
Oomycota use DAP; animals cannot make lysine
at all \cite[Chapter 7]{deacon2006fungal}.
Most fungi have ergosterol as their main
membrane sterol, which is why many
antifungal drugs (azoles, allylamines)
target its synthesis
\cite[Chapter 13]{webster2007fungi}.
Daniel's systematics essay adds class IV
chitin synthase to these markers.

\section{Fungus-like organisms}
\label{section-fungus-like}

Mycologists have traditionally also studied
organisms that look and live like fungi but
arose independently of the true fungi
(Eumycota). Webster and Weber place them in
two other kingdoms
\cite[Section 1.5]{webster2007fungi}:
\begin{itemize}
\item Protozoa: the slime moulds
(Myxomycota: plasmodial Myxomycetes and
cellular Dictyostelids) and the
plant-parasitic Plasmodiophoromycota.
\item Straminipila (Chromista):
Hyphochytriomycota, Labyrinthulomycota and,
above all, the Oomycota (water moulds,
downy mildews, \emph{Phytophthora}). The
group is named after the tinsel
(straminipilous) flagellum, and is related
to diatoms and brown algae.
\end{itemize}

Daniel's last exam outline lists why the
Oomycota are not fungi: cellulose (not
chitin) walls, lysine by the DAP pathway,
tubular mitochondrial cristae, a diploid
thallus, and zoospores with a
straminipilous flagellum
\cite[Chapter 4]{webster2007fungi}.
Their hyphae are coenocytic, as in
zygomycetes, which is why they were once
lumped with them as ``Phycomycetes''.
\emph{Phytophthora infestans}, the cause of
the Irish potato famine of the 1840s, is an
oomycete \cite[Chapter 1]{deacon2006fungal}.

\section{Thallus, spores and sexual cycle}
\label{section-thallus-spores}

\textbf{Thallus.}
The vegetative body may be
\cite[Section 1.1]{webster2007fungi}:
\begin{itemize}
\item a mycelium of hyphae, either
coenocytic (aseptate, many nuclei in one
cytoplasm; Oomycota, zygomycetes) or septate
(Ascomycota, Basidiomycota); septa are
usually perforated;
\item yeast cells, dividing by budding or
fission; dimorphic fungi switch between
yeast and hyphae;
\item a chytrid thallus, a walled body with
root-like rhizoids;
\item a naked plasmodium or a
pseudoplasmodium of amoeboid cells (in the
fungus-like groups).
\end{itemize}
A thallus is \emph{holocarpic} if all of it
turns into a reproductive structure, and
\emph{eucarpic} if reproductive organs form
on a part of it while the rest goes on
feeding.

\begin{definition}
\label{definition-dikaryon}
A mycelium is \emph{homokaryotic} if all its
nuclei are genetically identical and
\emph{heterokaryotic} if it contains nuclei
of different genotypes. It is
\emph{dikaryotic} if each cell contains two
genetically distinct haploid nuclei, one
from each parent; such a cell is written
$n+n$.
\end{definition}

\textbf{Spores.}
Asexual (mitotic) spores arise by
fragmentation of hyphae (arthrospores,
thick-walled chlamydospores), by budding, or
by sporulation: endogenously inside a
sporangium (motile planospores or zoospores;
non-motile aplanospores) or exogenously as
conidia on conidiophores. Sexual spores are
named after the process that makes them:
oospores (Oomycota), zygospores
(zygomycetes), ascospores (Ascomycota) and
basidiospores (Basidiomycota).

\textbf{The sexual cycle.}
Every sexual cycle has three steps:
\begin{enumerate}
\item \emph{plasmogamy}, the fusion of two
cytoplasms;
\item \emph{karyogamy}, the fusion of the
two nuclei into a diploid zygote nucleus;
\item \emph{meiosis}, which restores the
haploid state.
\end{enumerate}
In Ascomycota and Basidiomycota plasmogamy
and karyogamy are separated in time, and the
dikaryon lives in between, briefly in
ascomycetes and as the main feeding
mycelium in basidiomycetes. Plasmogamy may
happen between motile gametes (iso- or
anisogamy), by contact or conjugation of
gametangia, by spermatization (a small
spermatium fuses with a receptive female
cell), or by somatogamy (fusion of ordinary
hyphae). A species is \emph{homothallic} if
one individual can complete the sexual
cycle alone and \emph{heterothallic} if two
compatible mating types are needed.

\section{Classification}
\label{section-classification}

DNA sequences overturned the classical
classification. Hibbett et al.\ (2007)
recognize one kingdom Fungi with seven phyla
\cite{hibbett2007higher}:
\begin{itemize}
\item three phyla of flagellated fungi:
Chytridiomycota (in a narrow sense),
Neocallimastigomycota (anaerobic rumen
fungi) and Blastocladiomycota;
\item Microsporidia, intracellular
parasites;
\item Glomeromycota, the arbuscular
mycorrhizal fungi;
\item Ascomycota and Basidiomycota, joined
in the subkingdom Dikarya, named after the
dikaryon, their shared derived character.
\end{itemize}
The old Chytridiomycota and Zygomycota are
not monophyletic. Zygomycota is abandoned;
its members go to Glomeromycota and to four
subphyla of uncertain position
(Mucoromycotina, Entomophthoromycotina,
Kickxellomycotina, Zoopagomycotina).
Ascomycota has three subphyla:
Taphrinomycotina (an early lineage, e.g.\
\emph{Taphrina}, \emph{Schizosaccharomyces}),
Saccharomycotina (budding yeasts) and
Pezizomycotina (the filamentous ascomycetes
with ascomata). Basidiomycota also has
three: Pucciniomycotina (rusts),
Ustilaginomycotina (smuts) and
Agaricomycotina (mushrooms, brackets, jelly
fungi) \cite{hibbett2007higher}.

The course still used the older scheme of
Webster's and Alexopoulos's textbooks:
Zygomycota with classes Zygomycetes and
Trichomycetes; Ascomycota with
Archiascomycetes, Hemiascomycetes and
Euascomycetes; Basidiomycota with
Hymenomycetes, Gasteromycetes,
Urediniomycetes and Ustilaginomycetes
\cite[Section 1.5]{webster2007fungi}.

\textbf{The form group Deuteromycota.}
Fungi may have a sexual form (the
teleomorph, formerly the ``perfect state'')
and one or more asexual forms (anamorphs).
Fungi known only from their anamorph were
put in the artificial group Deuteromycota,
Fungi Imperfecti or mitosporic fungi
\cite[Section 1.4]{webster2007fungi}.
Culture work and DNA have tied most of them
to ascomycete (a few to basidiomycete)
teleomorphs; for example \emph{Botrytis
cinerea} is the anamorph of
\emph{Botryotinia fuckeliana}. Deuteromycetes
were classified by their conidia and
conidiomata: Hyphomycetes (free
conidiophores, or conidiophores bundled into
a synnema or a cushion-like sporodochium),
Coelomycetes (conidia inside a flask-shaped
pycnidium or a saucer-shaped acervulus under
the host epidermis) and Blastomycetes
(yeast-like anamorphs). Conidia are
\emph{thallic} if a pre-existing hyphal cell
becomes the spore (arthrospores,
chlamydospores) and \emph{blastic} if the
spore swells out from the conidiogenous cell
before a septum cuts it off (for example
phialides, which bud conidia from a fixed
opening).

\section{Life cycles of representative fungi}
\label{section-life-cycles}

\textbf{\emph{Rhizopus} and \emph{Mucor}
(Mucorales).}
The mycelium is coenocytic. \emph{Rhizopus}
spreads by stolons and anchors by rhizoids,
from which sporangiophores rise. Asexual
reproduction is by non-motile sporangiospores
formed in sporangia. In sex, two compatible
hyphae (mostly heterothallic) produce
gametangia that fuse (gametangial
conjugation); the fused cell becomes a thick-
walled zygosporangium, in which karyogamy
gives the resting zygospore; meiosis occurs
at germination
\cite[Chapter 7]{webster2007fungi}.
\emph{Pilobolus}, on herbivore dung, shoots
its whole sporangium towards the light so
that it lands on grass and is eaten again.

\textbf{\emph{Saccharomyces cerevisiae}
(Saccharomycotina).}
Both haploid and diploid cells reproduce by
budding. Cells of opposite mating type, a and
$\alpha$, fuse to a diploid;
on a poor medium the diploid cell turns
directly into an ascus, undergoes meiosis and
forms four ascospores, two of each mating
type. There is no ascoma
\cite[Chapter 10]{webster2007fungi}.

\textbf{\emph{Claviceps purpurea}
(ergot).}
An overwintered sclerotium (the ergot)
germinates into stalked stromata bearing
perithecia. Their ascospores, shot out at the
time the grasses flower, infect the stigma;
the mycelium grows into the ovary and forms
conidia in a sugary honeydew which insects
carry to other flowers. Finally the ovary is
replaced by a new sclerotium up to 3 cm long
\cite[Section 12.5]{webster2007fungi}.
Its alkaloids constrict blood vessels;
eating ergot-contaminated grain causes
ergotism, with gangrene of the limbs and
sometimes death
\cite[Section 12.5]{webster2007fungi}.

\textbf{\emph{Puccinia graminis} (black
stem rust).}
A macrocyclic and heteroecious rust: it has
all five spore stages and needs two hosts
\cite[Chapter 22]{webster2007fungi}.
On wheat, dikaryotic urediniospores (II)
spread the disease during summer, and
teliospores (III) form at the end of the
season; in the teliospore karyogamy occurs,
and on germination meiosis gives haploid
basidiospores (IV), which infect barberry
(\emph{Berberis}). There spermatia (0)
achieve plasmogamy between compatible
mycelia, and the resulting dikaryotic
aeciospores (I) can infect only wheat.

\textbf{\emph{Ustilago maydis} (corn
smut, huitlacoche).}
The fungus is dimorphic: a haploid,
saprotrophic yeast phase (sporidia) and a
dikaryotic parasitic mycelium. Two sporidia
fuse only if they differ at both mating-type
loci, $a$ and $b$; the $b$ locus controls
pathogenicity. The dikaryotic hyphae infect
maize and induce galls filled with dark
diploid teliospores; a teliospore germinates
into a promycelium (the basidium), where
meiosis yields new sporidia
\cite[Chapters 18 and 23]{webster2007fungi}.
The young galls are eaten in Mexico as
huitlacoche (cuitlacoche); Mexican farmers
welcome the infection because infected ears
fetch a higher price
\cite[pp.~254--255]{hall2003edible}.

\textbf{A mushroom (Agaricales).}
A basidiospore germinates into a
monokaryotic mycelium; two compatible
monokaryons fuse, and the dikaryon, which
keeps its two nuclei paired by clamp
connections at each septum, is the long-lived
feeding mycelium. The mushroom (basidioma) is
built of dikaryotic hyphae. Its gills or
tubes are lined by the hymenium of basidia,
sterile basidioles and cystidia. In each
basidium karyogamy and meiosis give
(usually) four basidiospores on sterigmata,
which are shot off a short distance
(ballistospores) and fall out between the
gills
\cite[Chapter 18]{webster2007fungi}.

\textbf{Ascomata and asci.}
In ascomycetes the female ascogonium
receives the male nucleus, often through a
receptive hair (trichogyne); dikaryotic
ascogenous hyphae then grow out, and the
asci form at their tips. Most asci contain
eight ascospores (meiosis followed by one
mitosis). Septa have a central pore that a
Woronin body can plug. The ascomata of
Daniel's outline are:
\begin{itemize}
\item cleistothecium: closed, spores freed
when the wall breaks (\emph{Aspergillus},
\emph{Penicillium}, powdery mildews);
\item perithecium: flask-shaped with an
apical pore, the ostiole (\emph{Neurospora},
\emph{Claviceps});
\item apothecium: an open cup or disc
(\emph{Peziza}, \emph{Sclerotinia}; truffles
are underground relatives);
\item ascostroma or pseudothecium: asci in
cavities dissolved in a stroma
(Loculoascomycetes).
\end{itemize}
Asci are prototunicate (the wall dissolves),
unitunicate (one functional wall, opening by
a lid, operculate, or by a pore,
inoperculate) or bitunicate (two walls; the
inner one extends through the burst outer
one)
\cite[Chapter 8]{webster2007fungi}.

\textbf{Basidiomata.}
The course used the old artificial groups:
agarics (gills), boletes (tubes), the
Aphyllophorales (brackets, coral fungi such
as \emph{Ramaria}, tooth fungi), the jelly
fungi with septate basidia
(phragmobasidia, in contrast with the
undivided holobasidia of mushrooms), and the
gasteromycetes (puffballs, earthstars such as
\emph{Geastrum}), whose spores ripen inside a
closed gleba and are not shot off.

\section{Fossils and divergence times}
\label{section-fossils}

Fossil fungi recognizable as chytrids,
zygomycetes and ascomycetes occur in the
Lower Devonian Rhynie chert, about 400
million years old, and the early land plants
there already had arbuscular mycorrhizae
\cite[Section 1.5]{webster2007fungi}.
Notable fossils \cite{taylor2006dating}:
\begin{itemize}
\item a glomeromycete from about 460 million
years ago;
\item \emph{Paleopyrenomycites devonicus},
from the Rhynie chert, a perithecial
ascomycete with ostiole, paraphyses and asci
of up to 32 ascospores.
\end{itemize}
Molecular clocks give very different dates
depending on calibration. Taylor and Berbee
date the Ascomycota--Basidiomycota split at
1808, 1489 or about 400 million years ago,
depending on whether they calibrate with the
bird--mammal split, with
\emph{Paleopyrenomycites} as a
sordariomycete, or with the
monocot--eudicot split
\cite{taylor2006dating}.
An earlier estimate put the separation of
fungi from animals near 900 million years
ago
\cite[Section 1.5]{webster2007fungi}.
Daniel's essay also discussed
\emph{Prototaxites}, a giant Devonian fossil
made of tubes, interpreted as a fungus or a
lichen.

\section{Ecology and importance}
\label{section-ecology}

\textbf{Decomposers.}
With bacteria, fungi are the main
decomposers; they are especially important
in breaking down cellulose, about 40\% of
plant cell walls; wood-rotting fungi such
as \emph{Phanerochaete chrysosporium} also
degrade lignin
\cite[Chapter 1]{deacon2006fungal}.

\textbf{Mycorrhizae.}
Mycorrhizae (``fungus roots'') occur in at
least 80\% of vascular plants: the fungus
passes mineral nutrients, especially
phosphorus, to the plant and receives
sugars. Arbuscular mycorrhizae, formed by
Glomeromycota inside root cells, are the
commonest type; ectomycorrhizae, formed
mostly by basidiomycetes on trees such as
pines and oaks, wrap the root in a sheath
and grow between the cortical cells (the
Hartig net)
\cite[Chapter 13]{deacon2006fungal}.

\textbf{Lichens.}
A lichen is a fungus (the mycobiont, usually
an ascomycete) with a photosynthetic partner
(the photobiont: a green alga or a
cyanobacterium). Lichens are named and
classified as fungi
\cite[Chapter 13]{deacon2006fungal}.

\textbf{Pathogens.}
Fungi and fungus-like organisms cause more
than 70\% of the major crop diseases.
Biotrophic parasites feed on living cells
(rusts, powdery mildews); necrotrophic
parasites kill host tissue as they feed
(\emph{Botrytis}). Only about 200 fungi
infect warm-blooded animals, mostly when the
immune system is weak
\cite[Chapter 1]{deacon2006fungal}.

\textbf{Daniel's \emph{Sclerotinia}
proposal.}
\emph{Sclerotinia sclerotiorum} is a
homothallic, necrotrophic ascomycete of the
Helotiales (inoperculate discomycetes)
\cite[Chapter 15]{webster2007fungi}.
Its sclerotia survive in soil and germinate
into small apothecia whose ascospores infect
the plant. In Manitoba in 1971 and 1975 more
than 90\% of sunflower plants were infected
with wilt and head rot
\cite[Section 4.4]{saharan2008sclerotinia}.
Daniel's course proposal was to model
sunflower (prey, $x$) and fungus (predator,
$y$) by the Lotka--Volterra equations
$$
\begin{aligned}
\frac{dx}{dt} &= \alpha x - \gamma x y, \\
\frac{dy}{dt} &= -\beta y + \varepsilon x y,
\end{aligned}
$$
with $\alpha$ the growth rate of the
sunflower without the fungus, $\beta$ the
death rate of the fungus without hosts, and
$\gamma$, $\varepsilon$ the effects of
encounters. The proposal stopped at the
model; it was not fitted to data.

\textbf{Food and industry.}
Fungi give bread, beer and wine
(\emph{Saccharomyces}), cheeses, edible
mushrooms and huitlacoche, antibiotics
(penicillins from \emph{Penicillium}),
ergot alkaloids, and industrial enzymes such
as amylases, proteases, pectinases and
lipases, largely from \emph{Aspergillus}
\cite[Chapter 1]{deacon2006fungal}.

\section{Mycetisms}
\label{section-mycetisms}

\begin{definition}
\label{definition-mycetism}
A \emph{mycetism} is poisoning by eating a
mushroom, that is, the fruiting body of a
fungus.
\end{definition}

About 100 known species of mushrooms are
poisonous to humans, and the list keeps
growing. Most poisonings follow foraging and
misidentification; most deaths are caused by
\emph{Amanita phalloides}
\cite{graeme2014mycetism}. The mycetism
project of Daniel's course used a pair of
look-alikes as its example: the edible
chanterelle \emph{Cantharellus cibarius} and
the toxic false chanterelle
\emph{Hygrophoropsis aurantiaca}. In
hospital, the syndrome (group of symptoms)
and its latency suggest the toxin, but
several genera may share one toxin
\cite[pp.~106--113]{hall2003edible}. Graeme
classifies poisonings by the organ system
most affected \cite{graeme2014mycetism}; the
table lists the main toxins.

\begin{center}
\begin{tabular}{ll}
toxin & main genera \\
\hline
amatoxins & \emph{Amanita},
\emph{Galerina}, \\
 & \emph{Lepiota}, \emph{Conocybe} \\
gyromitrin & \emph{Gyromitra} \\
orellanine & \emph{Cortinarius} \\
muscarine & \emph{Inocybe},
\emph{Clitocybe} \\
ibotenic acid, & \emph{Amanita muscaria},
\\
muscimol & \emph{A.\ pantherina} \\
psilocybin & \emph{Psilocybe},
\emph{Panaeolus} \\
coprine & \emph{Coprinus} \\
myotoxins & \emph{Tricholoma},
\emph{Russula} \\
\end{tabular}
\end{center}

\textbf{Amatoxins (liver).}
Amatoxins, chiefly $\alpha$-amanitin, are
heat-stable bicyclic octapeptides. They bind
RNA polymerase II irreversibly, stopping
mRNA and hence protein synthesis; cells that
take them up fastest (intestinal epithelium,
hepatocytes, kidney tubules) die
\cite{graeme2014mycetism}. Cooking does not
destroy them
\cite[pp.~106--107]{hall2003edible}.
The course follows four phases: a latency of
6 to 24 hours, a violent gastrointestinal
phase, an apparent remission while the liver
is being destroyed, and acute liver failure,
sometimes requiring a transplant
\cite{graeme2014mycetism}. Main species:
\emph{A.\ phalloides}, \emph{A.\ virosa},
\emph{A.\ verna}, \emph{A.\ bisporigera},
\emph{Galerina marginata} and small
\emph{Lepiota} species
\cite[pp.~106--107]{hall2003edible}. In
Mexico, where \emph{A.\ phalloides} is rare,
white amanitas cause the fatal cases; in one
documented case in Veracruz five of seven
people died after eating \emph{A.\ virosa}
\cite{perezmoreno1994multiple}. Gyromitrin
of the false morel \emph{Gyromitra
esculenta} also damages the liver; it breaks
down to toxic hydrazines
\cite[pp.~107--108]{hall2003edible}.

\textbf{Orellanine (kidney).}
Orellanine, a bipyridine N-dioxide of
\emph{Cortinarius orellanus} and \emph{C.\
speciosissimus}, concentrates in the kidney
and damages the tubular epithelium. Symptoms
appear after a latency of 36 hours to 21
days, so the cause is often missed; renal
failure may require dialysis or a
transplant. \emph{Amanita smithiana} causes
a different, faster renal syndrome
\cite{graeme2014mycetism}.

\textbf{Muscarine.}
Muscarine mimics acetylcholine at peripheral
muscarinic receptors, causing sweating,
salivation, tears, slow pulse and diarrhoea.
It is found in \emph{Inocybe} and small white
\emph{Clitocybe}; despite its name,
\emph{Amanita muscaria} contains very little
\cite{graeme2014mycetism}.

\textbf{Ibotenic acid and muscimol.}
Ibotenic acid and muscimol are derivatives of
glutamate and GABA. Ibotenic acid, an NMDA
glutamate receptor agonist, is unstable and
decarboxylates to muscimol, a GABA$_A$
receptor agonist. Symptoms begin 30 to 180
minutes after ingestion: excitation or
drowsiness, distorted perception,
hallucinations, ataxia, and coma only in
severe cases \cite{graeme2014mycetism}. The
main species are \emph{Amanita muscaria} and
\emph{A.\ pantherina}; recovery is usually
complete within a day and deaths are very
rare \cite{michelot2003amanita}.

\textbf{Psilocybin.}
Psilocybin is a prodrug: it is quickly
dephosphorylated to psilocin, which acts
mainly as an agonist of the serotonin
5-HT$_{2A}$ receptor. It makes
up 0.2 to 1\% of the dry weight of
hallucinogenic mushrooms (\emph{Psilocybe},
\emph{Panaeolus}, \emph{Conocybe},
\emph{Gymnopilus}, \emph{Pluteus}). Its
toxicity is low: the lethal dose is far above
the psychoactive one, and the deaths reported
come from accidents in the altered state
\cite{tyls2014psilocybin}.

\textbf{Gastrointestinal irritants and
coprine.}
Many species cause only vomiting and
diarrhoea within a few hours, for example
\emph{Agaricus xanthodermus}, \emph{Entoloma
sinuatum}, \emph{Omphalotus} and
\emph{Chlorophyllum molybdites}
\cite[pp.~110--113]{hall2003edible}.
\emph{Paxillus involutus} can in addition
cause an immune haemolytic anaemia after
years of eating it. Coprine of the inky cap
\emph{Coprinus atramentarius} inhibits
acetaldehyde dehydrogenase, so alcohol drunk
up to three days later causes a
disulfiram-like reaction
\cite{graeme2014mycetism}.

\textbf{Myotoxins.}
Eating \emph{Tricholoma equestre} at several
consecutive meals (three to nine meals, 100
to 400 g) can cause rhabdomyolysis, the
breakdown of muscle, sometimes reaching the
heart and killing; \emph{Russula
subnigricans} causes a similar syndrome
\cite{graeme2014mycetism}. Graeme also
describes sudden deaths in Yunnan, China,
linked to \emph{Trogia venenata}.

\section{Ethnomycology}
\label{section-ethnomycology}

This section is written from the textbooks;
the surviving coursework has no material on
it. (The market interviews of
Section~\ref{section-field-trip} are the one
exception.)

Michelot and Melendez-Howell describe two
traditional attitudes to mushrooms:
mycophilic peoples (Slavs, Lithuanians,
French, Italians, Spaniards) gather and eat
them, mycophobic ones (Germanic, Celtic,
British) avoid them. Languages keep a sharp
line between safe and suspect mushrooms:
``mushrooms'' and ``toadstools'' in English,
``setas'' and ``hongos venenosos'' in Spanish
\cite{michelot2003amanita}.

\textbf{The fly agaric.}
\emph{Amanita muscaria} is the emblematic
mushroom of tales and decorations. Its name
(``fly agaric'', ``matamoscas'', Fliegenpilz)
recurs across languages; the authors argue it
links the mushroom to madness and possession
rather than to a real insecticidal effect.
Shamans of several Siberian peoples (Ostyak,
Vogul, Kamchadal, Koryak, Chukchi) use it in
rites; since ibotenic acid and muscimol pass
into the urine, some drink the urine of
someone who ate it. R.~G. Wasson argued that
it is the \emph{soma} of the Rig Veda. In
Mexico it is eaten after peeling off the red
cuticle and discarding the cooking water,
which removes much of the water-soluble
toxins. The authors stress that the fly
agaric is not a true object of veneration,
in contrast with the \emph{Psilocybe}
mushrooms of Mesoamerica
\cite{michelot2003amanita}.

\textbf{Mexican sacred mushrooms.}
The ritual use of hallucinogenic mushrooms
in Mexico goes back some 3000 years and
continues today. Western science learned of
them in 1957 through Robert G. Wasson; Roger
Heim classified them, and Albert Hofmann
isolated and identified psilocybin in 1958
and synthesized it in 1959
\cite{tyls2014psilocybin}. The sources
opened for these notes do not discuss
Mar\'ia Sabina, the Mazatec healer
associated with this history.

\section{The field trip to Hidalgo}
\label{section-field-trip}

The course's collecting trip went to two
sites in the state of Hidalgo:
\begin{itemize}
\item the Parque Nacional El Chico, in the
mountains north of Pachuca, a cool, humid
forest of pine, oak and fir (\emph{Pinus},
\emph{Abies});
\item the reserve La Mojonera near
Zacualtip\'an de \'Angeles, a humid temperate
forest at 2040 m, with pine dominating lower
down and oak higher up.
\end{itemize}

\textbf{Method.}
For each specimen the team recorded date,
locality, vegetation and substrate;
photographed it in place; lifted it with a
trowel without damaging the fruiting body;
refilled the hole with substrate; wrapped it
in waxed paper or a paper bag; and labelled
it with a running number for description
back at camp, with a hand lens and a razor.

\textbf{What was found.}
All specimens were ascomycetes or
basidiomycetes, as expected for a collection
of fruiting bodies. At El Chico, on organic
matter: \emph{Gyromitra}, \emph{Peziza},
\emph{Russula}, \emph{Geastrum},
\emph{Agaricus}, yellow and white
\emph{Ramaria}, \emph{Marasmius} and
\emph{Clitocybe}; on wood: another
\emph{Peziza}, the jelly fungus
\emph{Dacrymyces} and a lichen. At La
Mojonera: the blue milk-cap \emph{Lactarius
indigo}, several \emph{Russula},
\emph{Clitocybe} and a tooth fungus
(Hydnaceae). Many specimens were spoiled by
moulds, among them an all-white
\emph{Amanita}, possibly one of the deadly
white species above. The hoped-for
\emph{Amanita muscaria}, recorded in the
area, was not found.

\textbf{The market.}
At the Zacualtip\'an market, two stall
holders sold bought mushrooms (oyster
mushrooms and huitlacoche) and knew little
about them. An elderly man who gathered his
own in Carpinteros sold ``manita'' (a
\emph{Ramaria}), ``ocotlapa'' and ``pancita
de venado'', and said that after many years
of gathering he could tell the poisonous
ones apart without trouble. The team
concluded that local knowledge of wild
mushrooms survives but that gathering for
sale is becoming rare.

\bibliography{my}
\bibliographystyle{amsalpha}

\end{document}
